\chapter{项目目标与目录结构}

\section{项目要解决的问题}

输入并不是一套已经可直接打印的双导结构，而是三个处于同一世界坐标系中的网格：
原始牙科导板、双导管装配体和患者牙列。程序需要保留导管装配关系，在导板上形成与
导管同轴的孔和足够的操作空间，再把导管、导板和新增梁架连接为一个可制造实体。

这项任务同时包含四类约束：

\begin{itemize}
  \item \term{功能约束}：导孔、C 口、操作窗和观察窗不能被融合重新堵住；
  \item \term{结构约束}：四根主梁、Y 梁及根部强化必须连续并与导板、导管可靠相交；
  \item \term{解剖约束}：梁和窗口的位置必须跟随牙弓语义，避免只依赖屏幕左右或固定 XYZ；
  \item \term{制造约束}：最终 STL 应闭合、流形、单一主要连通体，并控制孔道和净距。
\end{itemize}

\section{输入与 tooth-47 示例}

\begin{figure}[H]
  \centering
  \begin{subfigure}[t]{0.32\textwidth}
    \includegraphics[width=\linewidth]{01_input_template.png}
    \caption{原始牙科导板}
  \end{subfigure}\hfill
  \begin{subfigure}[t]{0.32\textwidth}
    \includegraphics[width=\linewidth]{02_input_sleeves.png}
    \caption{双导管装配体}
  \end{subfigure}\hfill
  \begin{subfigure}[t]{0.32\textwidth}
    \includegraphics[width=\linewidth]{03_input_dentition.png}
    \caption{患者牙列/牙龈表面}
  \end{subfigure}
  \caption{tooth-47 的三类核心输入。三者必须采用相同坐标、毫米单位，并在导出 STL 前应用对象变换。}
  \label{fig:inputs}
\end{figure}

tooth-47 还配置了手机 STL、手机止挡报告及一个旧实验 cutter。当前正式流程把手机
STL 和止挡报告用于运动包络；旧 cutter 并不是主流程的权威观察窗来源。

\section{顶层目录及职责}

\begin{longtable}{p{0.28\textwidth}p{0.62\textwidth}}
  \toprule
  目录或文件 & 主要职责\\
  \midrule
  \endfirsthead
  \toprule
  目录或文件 & 主要职责\\
  \midrule
  \endhead
  \projectpath{src/twin_guide/} & 核心包：配置、阶段规划、几何算法、Blender 建模、验证与公开 API。\\
  \projectpath{src/twin_guide/blender/} & STL 读写、网格查询、布尔运算、圆管扫掠、体素融合、渲染和导出。\\
  \projectpath{src/twin_guide/sleeve_estimation/} & 导管轴线拟合、切片、圆拟合、参数化重建与网格完整性检查。\\
  \projectpath{data/cases/} & 规范病例输入和病例 YAML；按单颗与多颗病例分开。\\
  \projectpath{examples/} & 可执行 JSON 工程配置，负责输入路径、工程参数和输出位置。\\
  \projectpath{output/} & 过程 JSON、调试网格、渲染图和最终 STL。\\
  \projectpath{docs/} & 使用指南、七阶段流程、架构、配置和 API 文档。\\
  \projectpath{tests/} & 纯几何单元测试、Blender 后端测试、病例端到端测试和批量回归入口。\\
  \projectpath{run.py} & Blender 内的命令行启动脚本。\\
  \projectpath{blender-env.sh} & 将项目源码、工作区模块和局部依赖注入 Blender 5.2 Python。\\
  \bottomrule
\end{longtable}

\section{配置分层}

TwinGuide 刻意把配置分成两层。JSON 是可执行工程配置，保存 STL 路径、导管尺寸、梁径、
体素精度、手机避让和输出目录；病例 YAML 保存上下颌、牙位集合、观察窗范围、锚点站位、
Y 梁模式和人工审核状态。读取时，\texttt{CaseConfig.from\_json()} 以 JSON 所在目录解析
相对路径，再加载 YAML 并合并设计字段。

这种分层的优点是工程参数和病例语义互不混淆；代价是必须维护二者的一致性。代码通过
拒绝未知字段、拒绝 YAML 重复键、拒绝 JSON/YAML 对同一设计字段重复定义等方式降低歧义。

\section{运行入口}

\begin{lstlisting}[language=bash,caption={当前项目的三个主要命令}]
./blender-env.sh --background --python run.py -- \
  process --config examples/case-tooth-47.json

./blender-env.sh --background --python run.py -- \
  generate --config examples/case-tooth-47.json --validate

./blender-env.sh --background --python run.py -- validate \
  --config examples/case-tooth-47.json \
  --model output/tooth_47/twin_guide.stl
\end{lstlisting}

\texttt{process} 只执行几何规划；\texttt{generate} 再把规划实体化并导出；
\texttt{validate} 读取最终 STL，重新建立预期几何并进行独立检查。生产生成还会检查 YAML
中的 review 状态，显式为 pending、unreviewed 等状态时默认停止。

